\clearpage
\subsection{联系两种构造}\label{sec:17g}
给定 $(c,k)$，可考虑 $\tilde d=\mu^2(c,\cdot)$ 及 $\tilde l=\mu^2(k(\cdot),\cdot)+\mu^3(c,\cdot,\cdot)$。它们满足\eref{17.15}，并可与 $(d,l)$ 联系。这可视为此前关于\eref{17.15}的解的唯一性断言的一个极限情形，但由于对应用很重要，我们采用独立记号，给出完整细致的论述。具体地，将构造
\begin{equation}\label{eq:17.19}
\boxed{\begin{gathered}
\xi:CF^{\mathrm{pr}}(Q,L_1)\longrightarrow CF^{\mathrm{pr}}(Q,L_2),\\
\psi:CF^{\mathrm{pr}}(L_0,L_1)\otimes CF^{\mathrm{pr}}(Q,L_0)\longrightarrow CF^{\mathrm{pr}}(Q,L_2),\\
\mu^1(\xi(\cdot))+\xi(\mu^1(\cdot))+d+\mu^2(c,\cdot)=0,\\
\mu^1(\psi(\cdot,\cdot))+\psi(\mu^1(\cdot),\cdot)+\psi(\cdot,\mu^1(\cdot))\\
{}+\xi(\mu^2(\cdot,\cdot))+l(\cdot,\cdot)
+\mu^2(k(\cdot),\cdot)+\mu^3(c,\cdot,\cdot)=0.
\end{gathered}}
\end{equation}
对于 $\xi$，先令 $E^{\xi,0}\to S^{\xi,0}$ 就是 $E^d\to S^d$，带之前采用的边界条件及相对扰动数据。在变形的另一端，取三标点圆盘上的平凡纤维化，三个边界分支上的边界条件为 $Q,L_1,L_2$，将其第二个出向端与 $E^c\to S^c$ 的入向端粘合。在两者间插值的族显而易见，见\sr{fig:17.7}{图~17.7}。
\SourceFigure{figure-17-7}{0.53169174}{17.7}

开始定义 $\psi$ 时，先写出五个单参数 Lefschetz 纤维化族，配备边界条件及扰动数据，
\SourcePage{260}
使它们合起来构成闭链：第一族的终点是第二族的起点，依此类推；图示见\sr{fig:17.8}{图~17.8}。第一族是之前定义的 $E^{l,r}\to S^{l,r}$。为了得到第二族，记为 $E^{\xi(\mu^2),r}\to S^{\xi(\mu^2),r}$，取 $E^{\xi,r}\to S^{\xi,r}$，将唯一的出向端与三标点圆盘上的平凡纤维化的入向端粘合；后者带边界条件 $Q,L_0,L_1$，以及定义 $\mu^2$ 乘积的扰动数据。因此，只要粘合长度充分大，计数该纤维化族的伪全纯截面，恰好给出 $\xi(\mu^2(\cdot,\cdot))$。回顾，$E^{l,1}\to S^{l,1}$ 是取 $E^d\to S^d$ 再粘入三标点圆盘上的平凡纤维化得到的；另一方面，由定义有 $(E^{\xi,0}\to S^{\xi,0})=(E^d\to S^d)$。这意味着，若适当选取粘合参数，则
\begin{equation}\label{eq:17.20}
(E^{\xi(\mu^2),0}\to S^{\xi(\mu^2),0})=(E^{l,1}\to S^{l,1}),
\end{equation}
其中等式也包括边界条件与扰动数据。在另一端 $r=1$，得到纤维化 $E^{\xi(\mu^2),1}\to S^{\xi(\mu^2),1}$，给出表达式 $\mu^2(c,\mu^2(\cdot,\cdot))$。这立即暗示下一步：在 $\RG^4$ 中取一个充分大的闭区间 $[0,1]$，考虑其上的四标点圆盘万有族，边界分支标记为 $(Q,L_0,L_1,L_2)$。\znote{252}{$\RG^4$ 是四个边界标点的圆盘模空间：按 (9.7)，$\RG^{d+1}=\operatorname{Conf}_{d+1}(\partial D)/\operatorname{Aut}(D)$，实维数为 $d-2$，故 $\RG^4$ 是一维区间，其 Deligne–Mumford 紧化把两端各缩成一点。所谓“充分大的闭区间”即伸到这两端附近，使某个标点被挤出的两种退化（粘合参数趋于零）也落在族内，于是这个单参数族的两端恰好与相邻的族拼接。[来源：Seidel, Fukaya categories and Picard–Lefschetz theory, §9 中 (9.7) 与 §12g。]}将其变成带适当边界条件与扰动数据的平凡 Lefschetz 纤维化族，再将每个圆盘的第三个出向端与 $E^c\to S^c$ 的入向端粘合。将结果记为 $E^{\mu^3(c\otimes\Id),r}\to S^{\mu^3(c\otimes\Id),r}$，因为计数该族的伪全纯截面就得到 $\mu^3(c,\cdot,\cdot)$。适当选择粘合参数，显然可使
\SourcePage{261}
\SourceFigure{figure-17-8}{0.60438396}{17.8}
\[
(E^{\mu^3(c\otimes\Id),0}\to S^{\mu^3(c\otimes\Id),0})
=(E^{\xi(\mu^2),1}\to S^{\xi(\mu^2),1}),
\]
等式的意义与\eref{17.20}相同。在另一端 $r=1$，得到的 Lefschetz 纤维化的截面定义同态 $\mu^2(\mu^2(c,\cdot),\cdot)$。对于第四族 $E^{\mu^2(k\otimes\Id),r}\to S^{\mu^2(k\otimes\Id),r}$，取边界分支标记为 $(Q,L_0,L_2)$ 的三标点圆盘，将其变成平凡 Lefschetz 纤维化，随后
\SourcePage{262}
将它的第二个出向端与纤维化 $E^{k,r}\to S^{k,r}$ 的入向端粘合。由 $k$ 的定义，
\[
(E^{\mu^2(k\otimes\Id),0}\to S^{\mu^2(k\otimes\Id),0})
=(E^{\mu^3(c\otimes\Id),1}\to S^{\mu^3(c\otimes\Id),1}).
\]
另一方面，在 $r=1$ 时，可将 $E^{\mu^2(k\otimes\Id),1}$ 视为如下构造：从三标点圆盘上的平凡纤维化出发，两个边界分支上的边界条件恒定等于 $Q,L_0$，余下分支上则给出从 $\tau_V(L_1)$ 到 $L_2$ 的同痕。在 $L_0$ 分支上标记一个边界点，如 $k$ 的定义那样，与带标准边界条件的 $E$ 作边界连通和。注意，$E^{l,0}\to S^{l,0}$ 的定义实质上相同。细节，例如标点位置，可能不同，但显然可用第五个也就是最后一个族在两者之间插值，记为 $E^{0,r}\to S^{0,r}$；其性质是，每个 $E^{0,r}$ 都由与 $E$ 作边界连通和得到，且粘合参数\emph{始终保持很小}。由定义，这使我们的链闭合：
\[
\begin{aligned}
(E^{0,0}\to S^{0,0})&=(E^{\mu^2(k\otimes\Id),1}\to S^{\mu^2(k\otimes\Id),1}),\\
(E^{0,1}\to S^{0,1})&=(E^{l,0}\to S^{l,0}).
\end{aligned}
\]
\stref{笔记}{17.8}的消失结果以适当的带参数版本适用于 $E^{0,r}\to S^{0,r}$，因此相应零维伪全纯截面模空间的计数平凡。最后，整个闭链形成以 $S^1$ 为参数的 Lefschetz 纤维化族，显然可用 $D^2$ 上的族填充。用所得带参数模空间定义 $\psi$，则\eref{17.19}由构造成立。

\StatementNumber{11}\begin{note}\label{st:17.11}
尝试将这一程序纳入更一般、更标准化的框架颇有启发。令 $\RG^{3,1}$ 为带一个额外内部标点的三标点圆盘模空间。等价地，也更适合眼前目的，可固定内部点及第一个边界点，从而消除 $\operatorname{Aut}(D)$ 对称性，并将 $\RG^{3,1}$ 看成另外两个边界点的构形空间。取圆盘 $D$ 为标准 Lefschetz 纤维化 $E\to S$ 的底空间，内部标点为唯一临界值，第一个边界点为 $*$。对任意构形 $(z_1,z_2)\in\RG^{3,1}$，可将 $E$ 限制到 $S\setminus\{*,z_1,z_2\}$，使该限制成为带条带状端的 Lefschetz 纤维化。为这一纤维化族赋予由 $Q,L_0,L_1$ 及从 $\tau_V(L_1)$ 到 $L_2$ 的同痕组成的 Lagrangian 边界条件，并考虑由所得带参数模空间给出的不变量。

为理解它所满足的关系，需要考察 Deligne–Mumford 型紧化 $\overline{\RG}^{3,1}$，它是一个六边形。与 $\overline{\RG}^5$ 一样，可将其实现为复 Deligne–Mumford 空间 $\overline{\M}_{0,5}$ 关于某个反全纯对合的实点集，
\SourcePage{263}
该对合将复共轭与两个标点的交换结合。对应于各层的退化示于\sr{fig:17.9}{图~17.9}。\znote{255}{六条边对应 $\overline{\RG}^{3,1}$ 的两类余维一退化：内部标点（即临界值）撞上某个边界标点，三种；两个边界标点相撞，也是三种；合计正是六，与图 17.9 周围六幅退化图一一对应。这也解释了何以只有边 (i)、(ii) 直接对应图 17.8 中的族：其余四条边给出的族只到“并的、保持端点不动的变形”为止，其贡献并入同一条链映射恒等式，而不成为独立项。[来源：Seidel, Fukaya categories and Picard–Lefschetz theory, §17g 与图 17.9。]}对于边 (i)、(ii)，所得单参数族是\sr{fig:17.8}{图~17.8}中相应族在粘合参数趋于零时的极限，给出\eref{17.19}中的项 $\mu^3(c\otimes\Id)$ 与 $\mu^2(k\otimes\Id)$。图的其余部分稍复杂：给出 $\xi(\mu^2)$ 项的族对应于边 (vi) 的一部分，$E^{0,r}\to S^{0,r}$ 也类似地覆盖边 (iii) 的一部分；这两条边的其余部分，以及 (iv)、(v)，在\sr{fig:17.8}{图~17.8}中没有直接对应物，但它们的并可保持端点不动地变形为给出 $l$ 的单参数族。因此，从图形上说，\sr{fig:17.8}{图~17.8}的五个族可看成\sr{fig:17.9}{图~17.9}中的虚线边。实际意义是：若直接使用 $\overline{\RG}^{3,1}$，则得到的是\eref{17.15}与\eref{17.19}稍复杂的版本，含有更多项。
\SourceFigure{figure-17-9}{0.60853781}{17.9}
\end{note}

\StatementNumber{12}\begin{note}\label{st:17.12}
在后续应用中，由于分次的存在，条件 $2c_1(M,V)=0$ 总会成立，见\secref{12a}的讨论。不过，必要时也可不用这一条件，而加入
\SourcePage{264}
涉及适当填充链的修正项，如\bcite{131}引理 3.12 或上述\stref{引理}{17.7}。
\end{note}
